\section{AIME I 2001}
        \begin{enumerate}
        	\item (\textit{AIME I 2001}) Find the sum of all positive two-digit integers that are divisible by each of their digits.


 
 \textbf{Answer: }630



\textbf{Solution 1}

Let our number be $10a + b$, $a,b \neq 0$. Then we have two conditions: $10a + b \equiv 10a \equiv 0 \pmod{b}$ and $10a + b \equiv b \pmod{a}$, or $a$ divides into $b$ and $b$ divides into $10a$. Thus $b = a, 2a,$ or $5a$ (note that if $b = 10a$, then $b$ would not be a digit).  

 
                \begin{itemize}
                    \item For $b = a$, we have $n = 11a$ for nine possibilities, giving us a sum of $11 \cdot \frac {9(10)}{2} = 495$.\item For $b = 2a$, we have $n = 12a$ for four possibilities (the higher ones give $b > 9$), giving us a sum of $12 \cdot \frac {4(5)}{2} = 120$.\item For $b = 5a$, we have $n = 15a$ for one possibility (again, higher ones give $b > 9$), giving us a sum of $15$.
                \end{itemize}


                If we ignore the case $b = 0$ as we have been doing so far, then the sum is $495 + 120 + 15 = \boxed{630}$.

 





\textbf{Solution 2}

By testing every 2-digit number, we can list out all of the possibilities: 
 \[11+12+15+22+24+33+36+44+48+55+66+77+88+99=\boxed{630}.\]

 





\textbf{Solution 3}

To further expand on solution 2, it would be tedious to test all $90$ two-digit numbers. We can reduce the amount to look at by focusing on the tens digit.First, we cannot have any number that is a multiple of $10$. We also note that any number with the same digits is a number that satisfies this problem. This gives 
 \[11, 22, 33, ... 99.\] We start from each of these numbers and constantly add the digit of the tens number of the respective number until we get a different tens digit. For example, we look at numbers $11, 12, 13, ... 19$ and numbers $22, 24, 26, 28$. This heavily reduces the numbers we need to check, as we can deduce that any number with a tens digit of $5$ or greater that does not have two of the same digits is not a valid number for this problem. This will give us the numbers from solution 2.

 





\textbf{Solution 4}

In this solution, we will do casework on the ones digit.Before we start, let's make some variables. Let $a$ be the ones digit, and $b$ be the 

 




	\item (\textit{AIME I 2001}) A finite set $\mathcal{S}$ of distinct real numbers has the following properties: the mean of $\mathcal{S}\cup\{1\}$ is $13$ less than the mean of $\mathcal{S}$, and the mean of $\mathcal{S}\cup\{2001\}$ is $27$ more than the mean of $\mathcal{S}$. Find the mean of $\mathcal{S}$.


 
 \textbf{Answer: }651



\textbf{Solution 1}

Let $x$ be the mean of $\mathcal{S}$. Let $a$ be the number of elements in $\mathcal{S}$.Then, the given tells us that $\frac{ax+1}{a+1}=x-13$ and $\frac{ax+2001}{a+1}=x+27$. Subtracting, we have 
 \begin{align*}\frac{ax+2001}{a+1}-40=\frac{ax+1}{a+1} \Longrightarrow \frac{2000}{a+1}=40 \Longrightarrow a=49\end{align*}
We plug that into our very first formula, and get:
 \begin{align*}\frac{49x+1}{50}&=x-13 \\ 49x+1&=50x-650 \\ x&=\boxed{651}.\end{align*}


 





\textbf{Solution 2}

Since this is a weighted average problem, the mean of $S$ is $\frac{13}{27}$ as far from $1$ as it is from $2001$. Thus, the mean of $S$ is$1 + \frac{13}{13 + 27}(2001 - 1) = \boxed{651}$.

 




	\item (\textit{AIME I 2001}) Find the sum of the roots, real and non-real, of the equation $x^{2001}+\left(\frac 12-x\right)^{2001}=0$, given that there are no multiple roots.


 
 \textbf{Answer: }500



\textbf{Solution 1}

By \href{https://artofproblemsolving.com/wiki/index.php?title=Vieta\%27s\_formulas}{Vieta's formulae}, for a \href{https://artofproblemsolving.com/wiki/index.php?title=Polynomial}{polynomial} equation of the form $a_nx^n + a_{n-1}x^{n-1} + \dotsb + a_0 = 0$, the sum of the roots, \textit{counted with multiplicity}, is $-\frac{a_{n-1}}{a_n}$. 

 Now, by the \href{https://artofproblemsolving.com/wiki/index.php?title=Binomial\_Theorem}{Binomial Theorem}, the highest-degree term of the expansion of $\left(\frac 12-x\right)^{2001}$ is $-x^{2001}$, but $x^{2001}+\left(-x^{2001}\right) = 0$, so the highest-degree term is in fact the $x^{2000}$ term. Since there are no multiple roots, counting them with multiplicity includes each root exactly once, so the required sum will be given by precisely the formula above.

 We accordingly need to compute the coefficients of the $x^{2000}$ and $x^{1999}$ terms, which is straightforward using the \href{https://artofproblemsolving.com/wiki/index.php?title=Binomial\_Theorem}{Binomial Theorem} again:

 
 \begin{align*}&\binom{2001}{1} \cdot (-x)^{2000} \cdot \left(\frac{1}{2}\right)^1 =\frac{2001}{2}x^{2000}, \quad\text{and} \\ &\binom{2001}{2} \cdot (-x)^{1999} \cdot \left(\frac{1}{2}\right)^2 =-\frac{2001 \cdot 2000}{8}x^{1999} = -2001 \cdot 250x^{1999}. \end{align*}


 Hence the sum of all the distinct roots is $-\frac{-2001 \cdot 250}{\left(\frac{2001}{2}\right)} = 250 \cdot 2 = \boxed{500}$.

 





\textbf{Solution 2}

As in Solution 1, we find that the $x^{2001}$ terms cancel, so the left-hand side of the given equation is a $2000^{\text{th}}$-degree polynomial. This means there are $2000$ roots, \textit{counted with multiplicity}, but as there are no multiple roots, all of these $2000$ roots are actually distinct. Moreover, noting the symmetry of the equation in the form given, we observe that if $x$ is a root, then $\left(\frac{1}{2}-x\right)$ is also a root. In each case, we must have $x \neq \frac{1}{2}-x$, since otherwise the solution of $x = \frac{1}{2}-x$, i.e. $x = \frac{1}{4}$, would be a multiple root. Thus we can pair up the $2000$ distinct roots to obtain $\frac{2000}{2} = 1000$ distinct pairs, each of which contains $2$ distinct roots that sum to $\frac{1}{2}$. Summing these pairs will therefore count each root exactly once, as desired, so the answer is $1000 \cdot \frac{1}{2} = \boxed{500}$.

 





\textbf{Solution 3}

As in Solution 2, observe that if $r$ is a root, then $\left(\frac{1}{2}-r\right)$ is also a root, yielding a pair of roots that sum to $\frac{1}{2}$. This motivates us to substitute $y = x-\frac{\left(\frac{1}{2}\right)}{2}$, i.e. $y = x-\frac{1}{4}$, so that the left-hand side of the equation becomes even more symmetric: 
 \[\left(\frac{1}{4}+y\right)^{2001}+\left(\frac{1}{4}-y\right)^{2001} = 0.\] Using the \href{https://artofproblemsolving.com/wiki/index.php?title=Binomial\_Theorem}{Binomial Theorem}, this expands as 
 \[2\cdot\frac{1}{4}\cdot\binom{2001}{1}y^{2000}-0y^{1999}+\dotsb = 0,\] so by \href{https://artofproblemsolving.com/wiki/index.php?title=Vieta\%27s\_formulas}{Vieta's formulae}, the sum of the roots for $y$ is 0. Since the original equation in $x$ has degree $2000$, there are also $2000$ roots for $x$, all of which are distinct since there are no multiple roots. Therefore, letting these roots be $x_n = y_n+\frac{1}{4}$ for $1 \leq n \leq 2000$, their sum is 
 \[\sum_{n=1}^{2000}\left(y_n+\frac{1}{4}\right) = \sum_{n=1}^{2000}y_n + 2000 \cdot \frac{1}{4} = 0 + 500 =\boxed{500}.\]

 





\textbf{Solution 4}

We begin by introducing a substitution. Let $x=\frac{1}{4}-y$. Substituting this into our equation gives 
 \[\left(\left(\frac{1}{4}\right)^{2001}-2001\left(\frac{1}{4}\right)^{2000}y+\binom{2001}{2}\left(\frac{1}{4}\right)^{1999}y^2+\cdots-y^{2001}\right)+\left(\left(\frac{1}{4}\right)^{2001}+2001\left(\frac{1}{4}\right)^{2000}y+\binom{2001}{2}\left(\frac{1}{4}\right)^{1999}y^2+\cdots+y^{2001}\right)=0.\] We can expand this out, canceling common terms, giving us, 
 \[2\left(\frac{1}{4}\right)^{2001}+2\binom{2001}{2}\left(\frac{1}{4}\right)^{1999}y^2+\cdots+2\binom{2001}{2000}\left(\frac{1}{4}\right)y^{2000}=0.\] We can clean it up by simply multiplying both sides by $2$, giving, 
 \[\frac{1}{4^{2000}}+\binom{2001}{2}\frac{y^2}{4^{1998}}+\binom{2001}{4}\frac{y^4}{4^{1996}}+\cdots+\binom{2001}{2000}y^{2000}.\] At first it may seem like we haven't progressed whatsoever, but we can revisit what the problem is asking. We aren't trying to find every root of this, but rather the sum of $\frac{1}{4}$ minus every root. Assume the roots were $r_1,r_2,...$. This means that our answer would be 

 




	\item (\textit{AIME I 2001}) In triangle $ABC$, angles $A$ and $B$ measure $60$ degrees and $45$ degrees, respectively. The bisector of angle $A$ intersects $\overline{BC}$ at $T$, and $AT=24$. The area of triangle $ABC$ can be written in the form $a+b\sqrt{c}$, where $a$, $b$, and $c$ are positive integers, and $c$ is not divisible by the square of any prime. Find $a+b+c$.


 
 \textbf{Answer: }291



\textbf{Solution}

After chasing angles, $\angle ATC=75^{\circ}$ and $\angle TCA=75^{\circ}$, meaning $\triangle TAC$ is an isosceles triangle and $AC=24$.

 Using law of sines on $\triangle ABC$, we can create the following equation:

 $\frac{24}{\sin(\angle ABC)}$ $=$ $\frac{BC}{\sin(\angle BAC)}$

 $\angle ABC=45^{\circ}$ and $\angle BAC=60^{\circ}$, so $BC = 12\sqrt{6}$.

 We can then use the Law of Sines area formula $\frac{1}{2} \cdot BC \cdot AC \cdot \sin(\angle BCA)$ to find the area of the triangle.

 $\sin(75)$ can be found through the sin addition formula.

 $\sin(75)$ $=$ $\frac{\sqrt{6} + \sqrt{2}}{4}$

 Therefore, the area of the triangle is $\frac{\sqrt{6} + \sqrt{2}}{4}$ $\cdot$ $24$ $\cdot$ $12\sqrt{6}$ $\cdot$ $\frac{1}{2}$

 $72\sqrt{3} + 216$

 $72 + 3 + 216 =$ $\boxed{291}$

 





\textbf{Solution 2 (no trig)}

First, draw a good diagram.

 We realize that $\angle C = 75^\circ$, and $\angle CAT = 30^\circ$.  Therefore, $\angle CTA = 75^\circ$ as well, making $\triangle CAT$ an isosceles triangle.  $AT$ and $AC$ are congruent, so $AC=24$. We now drop an altitude from $C$, and call the foot this altitude point $D$.  

 
                    \begin{figure}[H]
                        \centering
                        \includegraphics[width=0.5\textwidth]{Solution/AIME_I_2001_4_Solution_Figure_1.png}
                    \end{figure}
                    By 30-60-90 triangles, $AD=12$ and $CD=12\sqrt{3}$.  

 We also notice that $\triangle CDB$ is an isosceles right triangle.  $CD$ is congruent to $BD$, which makes $BD=12\sqrt{3}$.  The base $AB$ is $12+12\sqrt{3}$, and the altitude $CD=12\sqrt{3}$.  We can easily find that the area of triangle $ABC$ is $216+72\sqrt{3}$, so $a+b+c=\boxed{291}$.

 -youyanli

 





\textbf{Solution 3(Speedy and Simple)}

After drawing line $AT$, we see that we have two triangles: $\triangle ABT,$ with $45$, $30$, and $105$ degrees, and $\triangle ATC$, with $30$, $75$, $75$ degrees. If we can sum these two triangles' areas, we have our answer. 

 Let's take care of $\triangle ATC$ first. We see that $\triangle ATC$ is a isosceles triangle, with $AT = AC = 24$. Because the area of a triangle is $\frac{1}{2}ab\sin C$, we have $\frac{1}{2}\cdot 24^2\cdot\frac{1}{2}$, which is equal to $144.$ 

 Now, on to $\triangle ABT$. Draw the altitude from angle $\angle T$ to $AB$, and call the point of intersection $D$. This splits $\triangle ABT$ into $2$ triangles, one with $30-60-90$ ($\triangle ADT$), and another with $45-45-90$ ($\triangle BDT$). Now, because we know that $AT$ is $24$, we have by special right triangle ratios. The area of $\triangle ADT$ is $\frac{12\sqrt{3}\cdot 12}{2}$, and the area of $\triangle BDT$ is $\frac{12\cdot 12}{2}$, which adds to  $72\sqrt{3} + 72$. 

 Adding this to $\triangle ATC$ we get a total sum of $216 + 72\sqrt{3}.$ Thus, $a + b + c$ would be $216 + 72 + 3 = \boxed{291}.$

 





\textbf{Solution 4 (very fast)}

Recall the triangle area via sine formula $\frac{ab\sin{C}}{2}$. We notice that they have given almost all we need to use this, since $AC=24$ by properties of isosceles triangles and $\angle A$ itself equals $60$. So, we are trying to find $AB$. This is very trivial, as when we drop an altitude from $T$ to $AB$ (let the intersecting point be $U$), $AU=12\sqrt{3}$ and $BU=12$ by $30-60-90$ and $45-45-90$ triangles respectively. Therefore, the answer is just 
 \[\frac{(12+12\sqrt{3})(24)\sin{60}}{2}\]
 \[=(12+12\sqrt{3})(6)(\sqrt{3})\]
 \[=72\sqrt{3}+72\times 3\]
 \[=216 + 72\sqrt{3}\]
 \[\Longrightarrow 216+72+3=\boxed{291}.\]

 




	\item (\textit{AIME I 2001}) An equilateral triangle is inscribed in the ellipse whose equation is $x^2+4y^2=4$. One vertex of the triangle is $(0,1)$, one altitude is contained in the y-axis, and the length of each side is $\sqrt{\frac mn}$, where $m$ and $n$ are relatively prime positive integers. Find $m+n$.


 
 \textbf{Answer: }937



\textbf{Solution}


                    \begin{figure}[H]
                        \centering
                        \includegraphics[width=0.5\textwidth]{Solution/AIME_I_2001_5_Solution_Figure_1.png}
                    \end{figure}
                    





\textbf{Solution 1}

Denote the vertices of the triangle $A,B,$ and $C,$ where $B$ is in \href{https://artofproblemsolving.com/wiki/index.php?title=Quadrant}{quadrant} 4 and $C$ is in quadrant $3.$

 Note that the slope of $\overline{AC}$ is $\tan 60^\circ = \sqrt {3}.$ Hence, the equation of the line containing $\overline{AC}$ is
 \[y = x\sqrt {3} + 1.\]This will intersect the ellipse when
 \begin{eqnarray*}4 = x^{2} + 4y^{2} & = & x^{2} + 4(x\sqrt {3} + 1)^{2} \\ & = & x^{2} + 4(3x^{2} + 2x\sqrt {3} + 1) \implies x(13x+8\sqrt 3)=0\implies x = \frac { - 8\sqrt {3}}{13}. \end{eqnarray*}We ignore the $x=0$ solution because it is not in quadrant 3.

 Since the triangle is symmetric with respect to the y-axis, the coordinates of $B$ and $C$ are now $\left(\frac {8\sqrt {3}}{13},y_{0}\right)$ and $\left(\frac { - 8\sqrt {3}}{13},y_{0}\right),$ respectively, for some value of $y_{0}.$

 It is clear that the value of $y_{0}$ is irrelevant to the length of $BC$. Our answer is
 \[BC = 2 \cdot \frac {8\sqrt {3}}{13}=\sqrt {4\left(\frac {8\sqrt {3}}{13}\right)^{2}} = \sqrt {\frac {768}{169}}\implies m + n = \boxed{937}.\]

 





\textbf{Solution 2}

Solving for $y$ in terms of $x$ gives $y= -\frac {\sqrt{4-x^2}}{2}$, so the two other points of the triangle are $(x,-\frac {\sqrt{4-x^2}}{2})$ and $(-x,-\frac {\sqrt{4-x^2}}{2})$, which are a distance of $2x$ apart. Thus $2x$ equals the distance between $(x,-\frac {\sqrt{4-x^2}}{2})$ and $(0,1)$, so by the distance formula we have 

 
 \[2x=\sqrt{x^2+(\frac {2-\sqrt{4-x^2}}{2})^2}.\] 

 Squaring both sides and simplifying through algebra yields $x^2= \frac {192}{169}$, so $2x=\sqrt{\frac{768}{169}}$ and the answer is $\boxed{937}$.

 





\textbf{Solution 2 (slightly altered)}

Instead of using the distance between $(x,\frac {\sqrt{4-x^2}}{2})$ and $(0,1)$, we use the distance between the midpoint of the base and $(0,1)$.

 The midpoint of the base is $(0,-\frac {\sqrt{4-x^2}}{2})$, and the distance between the midpoint of the base and $(0,1)$ is equivalent to the height of the equilateral triangle: $x\sqrt{3}$.

 So,

 
 \[\frac {\sqrt{4-x^2}}{2}+1=x\sqrt{3}\]
 \[\frac {\sqrt{4-x^2}}{2}=x\sqrt{3}-1\]

 Squaring both sides:

 
 \[\frac {4-x^2}{4}=3x^2-2\sqrt{3}x+1\]
 \[4-x^2=12x^2-8\sqrt{3}x+4\]
 \[8\sqrt{3}x=13x^2\]
 \[x=\frac {8\sqrt{3}}{13}\]

 $x=\frac {8\sqrt{3}}{13}$, so $2x=\sqrt{\frac{768}{169}}$, and the answer is $768+169=\boxed{937}$.

 





\textbf{Solution 3}

Since the altitude goes along the $y$ axis, this means that the base is a horizontal line, which means that the endpoints of the base are $(x,y)$ and $(-x,y)$, and WLOG, we can say that $x$ is positive.

 Now, since all sides of an equilateral triangle are the same, we can do this (distance from one of the endpoints of the base to the vertex and the length of the base):

 $\sqrt{x^2 + (y-1)^2} = 2x$

 Square both sides,

 $x^2 + (y-1)^2 = 4x^2\implies (y-1)^2 = 3x^2$

 Now, with the equation of the ellipse:$x^2 + 4y^2 = 4$

 $x^2 = 4-4y^2$

 $3x^2 = 12-12y^2$

 Substituting, 

 $12-12y^2 = y^2 - 2y +1$

 Moving stuff around and solving:

 $y = \frac{-11}{13}, 1$

 The second is found to be extraneous, so, when we go back and figure out $x$ and then $2x$  (which is the side length), we find it to be:

 $\sqrt{\frac{768}{169}}$ 

 and so we get the desired answer of $\boxed{937}$.

 





\textbf{Solution 4}

Denote $(0,1)$ as vertex $A,$ $B$ as the vertex to the left of the $y$-axis and $C$ as the vertex to the right of the $y$-axis. Let $D$ be the intersection of $BC$ and the $y$-axis.

 Let $x_0$ be the $x$-coordinate of $C.$ This implies
 \[C=\left(x_0 , \sqrt{\frac{4-x_0^2}{4}}\right)\]and 
 \[B=\left(-x_0 , \sqrt{\frac{4-x_0^2}{4}}\right).\]Note that $BC=2x_0$ and 
 \[\frac{BC}{\sqrt3}=AD=1-\sqrt{\frac{4-x_0^2}{4}}.\]This yields
 \[\frac{2x_0}{\sqrt3}=1-\sqrt{\frac{4-x_0^2}{4}}.\]Re-arranging and squaring, we have
 \[\frac{4-x_0^2}{4}=\frac{4x_0^2}{3}-\frac{4x_0}{\sqrt3} +1.\]Simplifying and solving for $x_0$, we have
 \[x_0=\frac{48}{13\sqrt 3}.\]As the length of each side is $2x_0,$ our desired length is
 \[4x_0^2=\frac{768}{169}\]which means our desired answer is
 \[768+169=\boxed{937}\]

 





\textbf{Solution 5}


                    \begin{figure}[H]
                        \centering
                        \includegraphics[width=0.5\textwidth]{Solution/AIME_I_2001_5_Solution_Figure_2.png}
                    \end{figure}
                    


 We can begin by finding the length of line segment, $AE$ in terms of $m$ and $n$. On one hand, we can write, 
 \[AE= \frac{\sqrt{3}}{2}\cdot AC=\sqrt{\frac{3m}{4n}}\] as the altitude ($AE$) of an equilateral triangle is just $\frac{\sqrt{3}}{2}$ multiplied by length of a side ($AC$). One the other hand we can also write 
 \[AE=AO+OE\] since $A,O,E$ are colinear. Since $AO$ is the minor axis of the ellipse, it is equal to $1$. From the equation of the ellipse, we also get that 
 \[EC^2+4\cdot OE^2=4.\] Since 
 \[EC=\frac{1}{2} BC=\frac{1}{2}\sqrt{\frac{m}{n}}=\sqrt{\frac{m}{4n}}\] we get that 
 \[(\sqrt{\frac{m}{4n}})^2+4\cdot EO^2=4 \implies 4\cdot EO^2=4-\frac{m}{4n}\] so 
 \[EO=\sqrt{1-\frac{m}{16n}}.\] Substituting this value back into the equation, $AE=AO+OE$, we get that $AE=1+\sqrt{1-\frac{m}{16n}}$. Because both $AE=\sqrt{\frac{3m}{4n}}$, and $AE=1+\sqrt{1-\frac{m}{16n}}$, we see that 
 \[\sqrt{\frac{3m}{4n}}=1+\sqrt{1-\frac{m}{16n}}.\] Multiplying both sides by $4n$, we get, 
 \[\sqrt{12mn}=4n+\sqrt{16n^2-mn}.\] Thus, 
 \[\sqrt{12mn}-4n=\sqrt{16n^2-mn}.\] Hence, 
 \[(\sqrt{12mn}-4n)^2=(\sqrt{16n^2-mn})^2\] so 
 \[12mn+16n^2-8n\sqrt{12mn}=16n^2-mn.\] Canceling, we get 
 \[12mn-8n\sqrt{12mn}=-mn.\] We can rewrite this as 
 \[13mn=8n\sqrt{12mn}.\] Dividing both sides by $n$ gives us, $13m=8\sqrt{12mn}$. We can finally square both sides to get $169m^2=768mn$. Dividing both sides by $m$ gives $169m=768n$. Hence, 
 \[\frac{m}{n}=\frac{768}{169}.\] Since $m$ and $n$ are relatively prime, we see that $m=768$, and $n=169$. Therefore $m+n=768+169=\boxed{937}$

 


 





\textbf{Solution 6}

Notice that $x^2+4y^2=4$ can be rewritten as $(x)^2+(2y)^2=2^2$. The points of the triangle are $(0, 1)$, $(-x, 1-x\sqrt{3})$, and $(x, 1-x\sqrt{3})$. When plugging the second coordinate into the equation, we get $x^2+4-8x\sqrt{3}+12x^2=4$, which equals $13x^2-8x\sqrt{3}=0$. 

 This yields $x(13x-8\sqrt{3})=0$. Obviously x can't be 0, so $x=\frac{8\sqrt{3}}{13}$. The side length of the equilateral triangle is twice of this, so $\frac{16\sqrt{3}}{13}$.

 This can be rewritten as $\sqrt{\frac{256\cdot3}{169}}=\sqrt{\frac{768}{169}}$. 

 $768+169=\boxed{937}$.

 





\textbf{Solution 7}

Consider the transformation $(x,y)$ to $(x/2, y).$ This sends the ellipse to the unit circle. If we let $n$ be one-fourth of the side length of the triangle, the equilateral triangle is sent to an isosceles triangle with side lengths $2n, n\sqrt{13}, n\sqrt{13}.$ Let the triangle be $ABC$ such that $AB=AC.$ Let the foot of the altitude from A be $X.$ Then $BX=n,$ and $AX=2n\sqrt{3}.$ Let $C$ be a point such that $AC$ is a diameter of the unit circle. Then $XC=2-2n\sqrt{3}.$ Using power of a point on X,
 \[n^2=2n\sqrt{3}(2-2n\sqrt{3})\]Simplifying gets us to $13n^2=4n\sqrt{3}.$ Then, $n=\dfrac{4\sqrt{3}}{13}$ which means the side length is $\dfrac{16\sqrt{3}}{13}=\sqrt{\dfrac{768}{169}}.$Thus, the answer is $768+169=\boxed{937}.$

 




	\item (\textit{AIME I 2001}) A fair die is rolled four times. The probability that each of the final three rolls is at least as large as the roll preceding it may be expressed in the form $m/n$, where $m$ and $n$ are relatively prime positive integers. Find $m+n$.


 
 \textbf{Answer: }079



\textbf{Solution 1}

Recast the problem entirely as a block-walking problem.  Call the respective dice $a, b, c, d$.  In the diagram below, the lowest $y$-coordinate at each of $a$, $b$, $c$, and $d$ corresponds to the value of the roll.

 
                    \begin{figure}[H]
                        \centering
                        \includegraphics[width=0.5\textwidth]{Solution/AIME_I_2001_6_Solution_Figure_1.png}
                    \end{figure}
                    

 The red path corresponds to the sequence of rolls $2, 3, 5, 5$.  This establishes a \href{https://artofproblemsolving.com/wiki/index.php?title=Bijection}{bijection} between valid dice roll sequences and block walking paths.  

 The solution to this problem is therefore $\dfrac{\binom{9}{4}}{6^4} = {\dfrac{7}{72}}$. So the answer is $\boxed{079}$.

 





\textbf{Solution 1a (Very Similar But Explained Differently)}

Every die roll must be at least $1$. From there, you can make $5$ "increases" of value one until the highest possible value of a die ($6$), and once you "increase" a die roll a certain amount, all following die rolls are increased by that amount. Choose to distribute the indistinguishable "increases" among the $4$ die rolls or to never include an "increase" in any die roll, a total of $5$ distinguishable categories. By \href{https://artofproblemsolving.com/wiki/index.php?title=Ball-and-urn}{Ball-and-urn} there are $\binom{9}{4}$ outcomes that satisfy the problem. Continue with the solution above.

 





\textbf{Solution 2}

Let $a, b, c,$ and $d$ be the results of rolling the four dice respectively. We have the range $1\leq a\leq b\leq c\leq d\leq 6$, and there are $6^4=1296$ total outcomes from rolling the dice. To transfer the inequality into a strictly increasing inequality, we can transform it into $1\leq a<b+1<c+2<d+3\leq 9$. Now, lets suppose $a'=a, b'=b+1, c'=c+2,$ and $d'=d+3$. Then, $1\leq a'<b'<c'<d' \leq 9$. Clearly, there are $\binom{9}{4}=126$ values that satisfy this equation, so our answer is $\dfrac{\binom{9}{4}}{6^4} = {\dfrac{7}{72}}\implies \boxed{079}$.

 





\textbf{Solution 3}

If we take any \href{https://artofproblemsolving.com/wiki/index.php?title=Combination}{combination} of four numbers, there is only one way to order them in a non-decreasing order. It suffices now to find the number of combinations for four numbers from $\{1,2,3,4,5,6\}$. We can visualize this as taking the four dice and splitting them into 6 slots (each slot representing one of {1,2,3,4,5,6}), or dividing them amongst 5 separators. Thus, there are ${9\choose4} = 126$ outcomes of four dice. The solution is therefore $\frac{126}{6^4} = \frac{7}{72}$, and $7 + 72 = \boxed{079}$.

 





\textbf{Solution 4}

Call the dice rolls $a, b, c, d$. The difference between the $a$ and $d$ distinguishes the number of possible rolls there are. 

 
                \begin{itemize}
                    \item If $a - d = 0$, then the values of $b,\ c$ are set, and so there are $6$ values for $a,\ d$.\item If $a - d = 1$, then there are ${3\choose2} = 3$ ways to arrange for values of $b,\ c$, but only $5$ values for $a,\ d$.\item If $a - d = 2$, then there are ${4\choose2} = 6$ ways to arrange $b, c$, and there are only $6 - 2 = 4$ values for $a, d$.
                \end{itemize}


                Continuing, we see that the sum is equal to $\sum_{i = 0}^{5}{i+2\choose2}(6 - i) = 6 + 15 + 24 + 30 + 30 + 21 = 126$. The requested probability is $\frac{126}{6^4} = \frac{7}{72}$ and our answer is $\boxed{79}$.

 





\textbf{Solution 5}

The dice rolls can be in the form 

 ABCD
AABC
AABB
AAAB
AAAA

 where A, B, C, D are some possible value of the dice rolls. (These forms are not keeping track of whether or not the dice are in ascending order, just the possible outcomes.) 

 
            \begin{enumerate}
                \item Now, for the first case, there are ${6\choose4} = 15$ ways for this. We do not have to consider the order because the combination counts only one of the permutations; we can say that it counts the correct (ascending order) permutation.\item Second case: ${6\choose3} = 20$ ways to pick 3 numbers, ${3\choose1}$ ways to pick 1 of those 3 to duplicate. A total of 60 for this case.\item Third case: ${6\choose2}=15$ ways to pick 2 numbers. We will duplicate both, so nothing else in this case matters.\item Fourth case: ${6\choose2} = 15$ ways to pick 2 numbers. We pick one to duplicate with ${2\choose1} = 2$, so there are a total of 30 in this case.\item Fifth case: ${6\choose1} = 6$; all get duplicated so nothing else matters.
            \end{enumerate}
            There are a total of $6^4$ possible dice rolls. 

 Thus, 

 $\frac{m}{n} = \frac{15 + 60 + 15 + 30 + 6}{6^4} = \frac{126}{6^4} = \frac{7}{72}$Therefore, our answer is $\boxed{079}$

 





\textbf{Solution 6}

Consider the number of possible dice roll combinations which work after $1$ roll, after $2$ rolls, and so on. There is 6 possible rolls for the first dice. If the number rolled is a 1, then there are 6 further values that are possible for the second dice; if the number rolled is a 2, then there are 5 further values that are possible for the second dice, and so on. 

 Suppose we generalize this as a function, say $f(l,r)$ return the number of possible combinations after $l$ rolls and $r$ being the beginning value of the first roll. It becomes clear that from above, $f(1,r) = 1$; every value of $l$ after that is equal to the sum of the number of combinations of $l - 1$ rolls that have a starting value of at least $r$. If we slowly count through and add up all the possible combinations we get $\frac{7}{72}$ possibilities.

 





\textbf{Solution 7}

In a manner similar to the above solution, instead consider breaking it down into two sets of two dice rolls. The first \href{https://artofproblemsolving.com/wiki/index.php?title=Subset}{subset} must have a maximum value which is $\le$ the minimum value of the second subset. 

 
                \begin{itemize}
                    \item If the first subset ends in a 1, there is $1$ such subset and there are $6 + 5 + 4 + 3 + 2 + 1 = \frac{6}{2}(6 + 1) = 21$ ways of making the second subset.\item If the first subset ends in a 2, there is $2$ such subsets and there are $5 + \ldots + 1 = \frac{5}{2}(5 + 1) = 15$ ways of making the second subset.
                \end{itemize}


                Thus, the number of combinations is $\sum_{i=1}^6 i \cdot \left(\frac{(7 - i)(8 - i)}{2}\right) = 21 + 30 + 30 + 24 + 15 + 6 = 126$, and the probability again is $\frac7{72}$, giving $m+n=\boxed{079}$.

 





\textbf{Solution 8 Recursion Formula}

If you're too tired to think about any of the above smart transformations of the problem, a recursion formula can be a robust way to the correct answer. 

 We just need to work out the valid cases for each roll. Denote by $N_{k}(n)$ the number of valid cases in the $k+1$-th roll when the number $n$ is rolled, for $k=1,2,3$ and $1 \leq n \leq 6$. Then we have the following recursion formula:
 \[N_{1}(n) = n\]
 \[N_{2}(n) = N_{1}(1) + N_{1}(2) + ... + N_{1}(n)=N_{2}(n-1)+N_{1}(n)\]
 \[N_{3}(n) = N_{2}(1) + N_{2}(2) + ... + N_{2}(n)=N_{3}(n-1)+N_{2}(n)\]The logic is that, if $n$ is rolled, then the number of valid cases is the subtotal of valid cases in the preceding roll for the outcome of 1 to $n$.  The recursion can be easily calculated by hand when $N_k(n)$ are put in columns side by side, given the fact that te numbers are smaller than 100. Finally, the total number of cases is
 \[N = \sum_{n=1}^{6} N_{3}(n)\]and
 \[P = \frac{N}{6^4}\]

 





\textbf{Solution 9: Observation of each case}

Lets try casework and observe the cases. Notice that if the last roll is a $1$, then the only dice rolls may be $1-1-1-1$, which is only $1$ possibility. Observe that if the last roll is $2$, then there are $4 = 1 + 3$ possibilities. When the last roll is a $3$, there are $10 = 1 + 3 + 6$ possibilities. Notice when the last roll is $n$, the number of cases is the sum of the first $n$ positive triangular numbers. This is easily provable when observing the numbers of possibilities after assigning a value to the last and second to last rolls.So there are a total of $1 + 1 + 3 + 1 + 3 + 6 + 1 + 3 + 6 + 10 + 1 + 3 + 6 + 10 + 15 + 1 + 3 + 6 + 10 + 15 + 21 = 126$ possibilities. So the probability is $\frac{126}{6^4} = \frac{7}{72}$ and $7 + 72 = \boxed{079}$. ~skyscraper

 





\textbf{Solution 10: Distributions}

This is equivalent to picking a four-element sequence of $\{1, 2, 3, 4, 5, 6\}$ with repetition. Notice that once the sequence is picked, there is one and only one way to order these so that they form a sequence of rolls satisfying our conditions. 

 Now count the number of such four-element sequences, let $a$ be the number of $1$s in the sequence, $b$ be the number of $2$s, $c$ $3$s, $d$ $4$s, $e$ $5$s, and $f$ $6$s. Now we see that we must have 
 \[a + b + c + d + e + f = 4\] with $a, b, c, d, e, f$ being nonnegative integers since there are a total of $4$ numbers picked. The number of solutions to this is $\dbinom{9}{4},$ so our total number is equal to $\dfrac{\binom{9}{4}}{6^4} = \dfrac{7}{72},$ making our answer $\boxed{079}.$ 

 





\textbf{Solution 11: Stars and Bars (Under 2 minutes)}

Let the rolls be $a,b,c$ and $d\newline$ let $z=a-1, e=b-a, f=c-b, g=d-c, h=6-d\newline$ $z+e+f+g+h=5\newline$ This equation has $C(5+5-1, 5-1)=126$ integer solutions$\newline$ $\frac{126}{1296}=\frac{7}{72}\newline$ $7+72=\boxed{079}$~ryanbear

 





\textbf{Solution 12}

Let's say the four dice values are all different. These can only be arranged in one way to satisfy our conditions, so there are $\binom{6}{4}=15$ ways. If there are three different values, there are $\binom{6}{3}$ to choose the numbers and $\binom{3}{1}$ to choose which number will have the repeat, so $20\cdot3=60$ ways. If there are two different values, there are $\binom{6}{2}$ ways to choose the numbers. If there are two of each, then there is one way. If there are three of one and one of another, then there are $\binom{2}{1}$ ways. Therefore $15\cdot(1+2)=45$. Now if all the numbers are the same, there are $\binom{6}{1}=6$ ways. Altogether we have $15+60+45+6=126$ ways. $\frac{126}{6^4}=\frac{21}{216}=\frac{7}{72}$. $7+72=\boxed{079}$. ~MC413551

 




	\item (\textit{AIME I 2001}) Triangle $ABC$ has $AB=21$, $AC=22$ and $BC=20$. Points $D$ and $E$ are located on $\overline{AB}$ and $\overline{AC}$, respectively, such that $\overline{DE}$ is parallel to $\overline{BC}$ and contains the center of the inscribed circle of triangle $ABC$. Then $DE=\frac{m}{n}$, where $m$ and $n$ are relatively prime positive integers. Find $m+n$.


 
 \textbf{Answer: }923



\textbf{Solution 1}


                    \begin{figure}[H]
                        \centering
                        \includegraphics[width=0.5\textwidth]{Solution/AIME_I_2001_7_Solution_Figure_1.png}
                    \end{figure}
                    Let $I$ be the \href{https://artofproblemsolving.com/wiki/index.php?title=Incenter}{incenter} of $\triangle ABC$, so that $BI$ and $CI$ are \href{https://artofproblemsolving.com/wiki/index.php?title=Angle\_bisector}{angle bisectors} of $\angle ABC$ and $\angle ACB$ respectively. Then, $\angle BID = \angle CBI = \angle DBI,$ so $\triangle BDI$ is \href{https://artofproblemsolving.com/wiki/index.php?title=Isosceles\_triangle}{isosceles}, and similarly $\triangle CEI$ is isosceles. It follows that $DE = DB + EC$, so the perimeter of $\triangle ADE$ is $AD + AE + DE = AB + AC = 43$. Hence, the ratio of the perimeters of $\triangle ADE$ and $\triangle ABC$ is $\frac{43}{63}$, which is the scale factor between the two similar triangles, and thus $DE = \frac{43}{63} \times 20 = \frac{860}{63}$. Thus, $m + n = \boxed{923}$.

 





\textbf{Solution 2}


                    \begin{figure}[H]
                        \centering
                        \includegraphics[width=0.5\textwidth]{Solution/AIME_I_2001_7_Solution_Figure_2.png}
                    \end{figure}
                    The \href{https://artofproblemsolving.com/wiki/index.php?title=Semiperimeter}{semiperimeter} of $ABC$ is $s = \frac{20 + 21 + 22}{2} = \frac{63}{2}$. By \href{https://artofproblemsolving.com/wiki/index.php?title=Heron\%27s\_formula}{Heron's formula}, the area of the whole triangle is $A = \sqrt{s(s-a)(s-b)(s-c)} = \frac{21\sqrt{1311}}{4}$. Using the formula $A = rs$, we find that the \href{https://artofproblemsolving.com/wiki/index.php?title=Inradius}{inradius} is $r = \frac{A}{s} = \frac{\sqrt{1311}}6$. Since $\triangle ADE \sim \triangle ABC$, the ratio of the heights of triangles $ADE$ and $ABC$ is equal to the ratio between sides $DE$ and $BC$. From $A=\frac{1}{2}bh$, we find $h_{ABC} = \frac{21\sqrt{1311}}{40}$. Thus, we have 

 $\frac{h_{ADE}}{h_{ABC}} = \frac{h_{ABC}-r}{h_{ABC}} = \frac{21\sqrt{1311}/40-\sqrt{1311}/6}{21\sqrt{1311}/40}=\frac{DE}{20}.$ Solving for $DE$ gives $DE=\frac{860}{63},$ so the answer is $m+n=\boxed{923}$.

 


 Or we have the area of the triangle as $S$.Using the ratio of heights to ratio of bases of $ADE$ and $ABC$$\frac {\frac {2S}{20}-\frac {2S}{63}}{\frac {2S}{20}}= \frac {DE}{BC(20)}$from that it is easy to deduce that $DE=\frac{860}{63}$.

 





\textbf{Solution 3 (mass points)}


                    \begin{figure}[H]
                        \centering
                        \includegraphics[width=0.5\textwidth]{Solution/AIME_I_2001_7_Solution_Figure_3.png}
                    \end{figure}
                    Let $P$ be the \href{https://artofproblemsolving.com/wiki/index.php?title=Incenter}{incenter}; then it is be the intersection of all three \href{https://artofproblemsolving.com/wiki/index.php?title=Angle\_bisector}{angle bisectors}. Draw the bisector $AP$ to where it intersects $BC$, and name the intersection $F$. 

 Using the \href{https://artofproblemsolving.com/wiki/index.php?title=Angle\_bisector\_theorem}{angle bisector theorem}, we know the ratio $BF:CF$ is $21:22$, thus we shall assign a weight of $22$ to point $B$ and a weight of $21$ to point $C$, giving $F$ a weight of $43$. In the same manner, using another bisector, we find that $A$ has a weight of $20$. So, now we know $P$ has a weight of $63$, and the ratio of $FP:PA$ is $20:43$. Therefore, the smaller similar triangle $ADE$ is $43/63$ the height of the original triangle $ABC$. So, $DE$ is $43/63$ the size of $BC$. Multiplying this ratio by the length of $BC$, we find $DE$ is $860/63 = m/n$. Therefore, $m+n=\boxed{923}$.

 





\textbf{Solution 4 (Faster)}

More directly than Solution 2, we have 
 \[DE=BC\left(\frac{h_a-r}{h_a}\right)=20\left(1-\frac{r}{\frac{[ABC]}{\frac{BC}{2}}}\right)=20\left(1-\frac{10r}{sr}\right)=20\left(1-\frac{10}{\frac{63}{2}}\right)=\frac{860}{63}\implies \boxed{923}.\]

 


 





\textbf{Solution 5}

Diagram borrowed from Solution 3.

 
                    \begin{figure}[H]
                        \centering
                        \includegraphics[width=0.5\textwidth]{Solution/AIME_I_2001_7_Solution_Figure_4.png}
                    \end{figure}
                    Let the angle bisector of $\angle{A}$ intersects $BC$ at $F$.

 Applying the \href{https://artofproblemsolving.com/wiki/index.php?title=Angle\_Bisector\_Theorem}{Angle Bisector Theorem} on $\angle{A}$ we have
 \[\frac{AB}{BF}=\frac{AC}{CF}\]
 \[BF=BC\cdot(\frac{AB}{AB+AC})\]
 \[BF=\frac{420}{43}\]Since $BP$ is the angle bisector of $\angle{B}$, we can once again apply the Angle Bisector Theorem on $\angle{B}$ which gives
 \[\frac{BA}{AP}=\frac{BF}{FP}\]
 \[\frac{AP}{PF}=\frac{AB}{BF}=\frac{41}{20}\]Since $\bigtriangleup ADE\sim\bigtriangleup ABC$ we have
 \[\frac{DE}{BC}=\frac{AP}{AF}\]
 \[DE=BC\cdot(\frac{AP}{(\frac{61}{41})\cdot AP})\]Solving gets $DE=\frac{860}{63}$. Thus $m+n=860+63=\boxed{923}$.

 





\textbf{Solution 6}

Let $A'$ be the foot of the altitude from $A$ to $\overline {BC}$ and $K$ be the foot of the altitude from $A$ to $\overline{DE}$. Evidently, 
 \[\frac{AK}{AA'} = 1- \frac{r}{AA'} = 1 - \frac{T/s}{T/BC}\] where $r$ is the inradius, $T = [ABC]$, and $s$ is the semiperimeter. So, 
 \[\frac{AK}{AA'} = 1 - \frac{BC}{s} = 1 - \frac{20}{63}= \frac{43}{63}\] Therefore, by similar triangles, we have $DE = BC * \frac{AK}{AA'} = 20 * \frac{AK}{AA'}= \boxed{\frac{860}{63}}$.

 


 





\textbf{Solution 7}

Label $P$ the point the angle bisector of $A$ intersects ${BC}$. First we find ${BP}$ and ${PC}$. By the Angle Bisector Theorem, $\frac{BP}{PC} = \frac{21}{22}$ and solving for each using the fact that ${BC} = 20$, we see that ${BP} = \frac{420}{43}$ and $PC = \frac{440}{43}$. 

 \begin{verbatim}
Because ${AP}$ is the angle bisector of $<A$, we can simply calculate it using Stewarts,
\end{verbatim}
 \[{AP} = 21*22 - \frac{440}{43}\cdot\frac{420}{43}\]
 \[{AP} = 21*22 - \frac{440\cdot420}{43^2}\]

 Now we can calculate what ${AO}$ is. Using the formula to find the distance from a vertex to the incenter, ${AO} = \frac{43}{63} \cdot[21\cdot22 - \frac{420*440}{43^2}] = \frac{43^2\cdot22 - 20\cdot440}{43\cdot3}$.

 Now because $\triangle{APE} ~ \triangle{ABC}$, we can find ${DE}$ by $\frac{AO}{AP} \cdot 20$. Dividing and simplifying, we see that $\frac{1}{21}\cdot\frac{43}{3}\cdot20 = \frac{860}{63}$. So the answer is $\boxed{923}$

 





\textbf{Solution 8 (vectors)}

To solve this problem, we can use the fact that, in $\triangle ABC$, the vector representation of the incenter is $\overrightarrow I = \frac{a\overrightarrow A + b\overrightarrow B + c\overrightarrow C}{a+b+c}$ and that that the vector of the foot of the bisector of $\angle BAC$ on $\overline{BC}$ is $\overrightarrow P = \frac{b\overrightarrow B + c\overrightarrow C}{b+c}$, where $a=BC,$ $b=AC,$ and $c=AB$.

 Let point $A$ be the origin of the coordinate plane. Then, $\overrightarrow A$ is the zero vector, so we can simplify our expression for $\overrightarrow I$ to $\frac{b\overrightarrow B + c\overrightarrow C}{a+b+c}$. Now, note that the vector components of $\overrightarrow I$ and $\overrightarrow P$ are the same, but they are multiplied by different scalars. Thus, the ratio of these scalars is the ratio of these vectors' magnitudes. Thus, we have $\frac{|\overrightarrow I|}{|\overrightarrow P|}=\frac{\frac1{a+b+c}}{\frac1{b+c}}=\frac{b+c}{a+b+c}=\frac{43}{63}$.

 Let $D \in \overline{AB}$ and $E \in \overline{AC}$. Because $\triangle AIE \sim \triangle APC$, we have $\frac{AI}{AP}=\frac{AE}{AC}$. Further, because $\triangle ADE \sim \triangle ABC$, we have $\frac{AE}{AC}=\frac{DE}{BC}$. Thus, by transitivity, $\frac{AI}{AP}=\frac{DE}{BC}$. We know that $\frac{AI}{AP}=\frac{43}{63}$, so $DE=\frac{AI}{AD}\cdot BC = \frac{43}{63}\cdot 20 = \frac{860}{63}$.

 Thus, our answer is $860+63=\boxed{923}$.

 




	\item (\textit{AIME I 2001}) Call a positive integer $N$ a $\textit{7-10 double}$ if the digits of the base-7 representation of $N$ form a base-10 number that is twice $N$. For example, $51$ is a 7-10 double because its base-7 representation is $102$. What is the largest 7-10 double?


 
 \textbf{Answer: }315



\textbf{Solution}

We let $N_7 = \overline{a_na_{n-1}\cdots a_0}_7$; we are given that

 
 \[2(a_na_{n-1}\cdots a_0)_7 = (a_na_{n-1}\cdots a_0)_{10}\] (This is because the digits in $N$ ' s base 7 representation make a number with the same digits in base 10 when multiplied by 2)

 Expanding, we find that

 
 \[2 \cdot 7^n a_n + 2 \cdot 7^{n-1} a_{n-1} + \cdots + 2a_0 = 10^na_n + 10^{n-1}a_{n-1} + \cdots + a_0\]

 or re-arranging,

 
 \[a_0 + 4a_1 = 2a_2 + 314a_3 + \cdots + (10^n - 2 \cdot 7^n)a_n\]

 Since the $a_i$s are base-$7$ digits, it follows that $a_i < 7$, and the LHS is less than or equal to $30$. Hence our number can have at most $3$ digits in base-$7$. Letting $a_2 = 6$, we find that $630_7 = \boxed{315}_{10}$ is our largest 7-10 double.

 





\textbf{Solution 2 (Guess and Check)}

Let $A$ be the base $10$ representation of our number, and let $B$ be its base $7$ representation.

 Given this is an AIME problem, $A<1000$. If we look at $B$ in base $10$, it must be equal to $2A$, so $B<2000$ when $B$ is looked at in base $10.$

 If $B$ in base $10$ is less than $2000$, then $B$ as a number in base $7$ must be less than $2*7^3=686$.

 $686$ is non-existent in base $7$, so we're gonna have to bump that down to $666_7$. 

 This suggests that $A$ is less than $\frac{666}{2}=333$. 

 Guess and check shows that $A<320$, and checking values in that range produces $\boxed{315}$.

 


 





\textbf{Solution 3}

Since this is an AIME problem, the maximum number of digits the 7-10 double can have is 3.  Let the number be 
 \[abc\] in base 7. Then the number in expanded form is 
 \[49a+7b+c\] in base 7 and 
 \[100a+10b+c\] in base 10.  Since the number in base 7 is half the number in base 10, we get the following equation.
 \[98a+14b+2c=100a+10b+c\] which simplifies to 
 \[2a=4b+c.\]  The largest possible value of a is 6 because the number is in base 7.  Then to maximize the number, $b$ is $3$ and $c$ is $0$.  Therefore, the largest 7-10 double is 630 in base 7, or $\boxed{315}$ in base 10.

 




	\item (\textit{AIME I 2001}) In triangle $ABC$, $AB=13$, $BC=15$ and $CA=17$. Point $D$ is on $\overline{AB}$, $E$ is on $\overline{BC}$, and $F$ is on $\overline{CA}$. Let $AD=p\cdot AB$, $BE=q\cdot BC$, and $CF=r\cdot CA$, where $p$, $q$, and $r$ are positive and satisfy $p+q+r=2/3$ and $p^2+q^2+r^2=2/5$. The ratio of the area of triangle $DEF$ to the area of triangle $ABC$ can be written in the form $m/n$, where $m$ and $n$ are relatively prime positive integers. Find $m+n$.


 
 \textbf{Answer: }061



\textbf{Solution 1}


                    \begin{figure}[H]
                        \centering
                        \includegraphics[width=0.5\textwidth]{Solution/AIME_I_2001_9_Solution_Figure_1.png}
                    \end{figure}
                    We let $[\ldots]$ denote area; then the desired value is 

 $\frac mn = \frac{[DEF]}{[ABC]} = \frac{[ABC] - [ADF] - [BDE] - [CEF]}{[ABC]}$Using the \href{https://artofproblemsolving.com/wiki/index.php?title=Triangle\#Related\_Formulae\_and\_Theorems}{formula} for the area of a triangle $\frac{1}{2}ab\sin C$, we find that

 $\frac{[ADF]}{[ABC]} = \frac{\frac 12 \cdot p \cdot AB \cdot (1-r) \cdot AC \cdot \sin \angle CAB}{\frac 12 \cdot AB \cdot AC \cdot \sin \angle CAB} = p(1-r)$and similarly that $\frac{[BDE]}{[ABC]} = q(1-p)$ and $\frac{[CEF]}{[ABC]} = r(1-q)$. Thus, we wish to find 
 \begin{align*}\frac{[DEF]}{[ABC]} &= 1 - \frac{[ADF]}{[ABC]} - \frac{[BDE]}{[ABC]} - \frac{[CEF]}{[ABC]}  \\ &= 1 - p(1-r) - q(1-p) - r(1-q)\\ &= (pq + qr + rp) - (p + q + r) + 1 \end{align*}
We know that $p + q + r = \frac 23$, and also that $(p+q+r)^2 = p^2 + q^2 + r^2 + 2(pq + qr + rp) \Longleftrightarrow pq + qr + rp = \frac{\left(\frac 23\right)^2 - \frac 25}{2} = \frac{1}{45}$. Substituting, the answer is $\frac 1{45} - \frac 23 + 1 = \frac{16}{45}$, and $m+n = \boxed{061}$.

 





\textbf{Solution 2}

By the barycentric area formula, our desired ratio is equal to 
 \begin{align*} \begin{vmatrix} 1-p & p & 0 \\  0 & 1-q & q  \\  r & 0 & 1-r \notag \end{vmatrix} &=1-p-q-r+pq+qr+pr\\ &=1-(p+q+r)+\frac{(p+q+r)^2-(p^2+q^2+r^2)}{2}\\ &=1-\frac{2}{3}+\frac{\frac{4}{9}-\frac{2}{5}}{2}\\ &=\frac{16}{45}, \end{align*}
 so the answer is $\boxed{061}$

 





\textbf{Solution 3 (Informal)}

Since the only conditions are that $p + q + r = \frac{2}{3}$ and $p^2 + q^2 + r^2 = \frac{2}{5}$, we can simply let one of the variables be equal to 0. In this case, let $p = 0$. Then, $q + r = \frac{2}{3}$ and $q^2 + r^2$ = $\frac{2}{5}$. Note that the ratio between the area of $DEF$ and $ABC$ is equivalent to $(1-q)(1-r)$. Solving this system of equations, we get $q = \frac{1}{3} \pm \sqrt{\frac{4}{45}}$, and $r = \frac{1}{3} \mp \sqrt{\frac{4}{45}}$. Plugging back into $(1-q)(1-r)$, we get $\frac{16}{45}$, so the answer is $\boxed{061}$

 \textbf{Note} Because the givens in the problem statement are all regarding the ratios of the sides, the side lengths of triangle $ABC$, namely $13, 15, 17$, are actually not necessary to solve the problem. This is clearly demonstrated in all of the above solutions, as the side lengths are not used at all.

 




	\item (\textit{AIME I 2001}) Let $S$ be the set of points whose coordinates $x,$ $y,$ and $z$ are integers that satisfy $0\le x\le2,$ $0\le y\le3,$ and $0\le z\le4.$  Two distinct points are randomly chosen from $S.$  The probability that the midpoint of the segment they determine also belongs to $S$ is $m/n,$ where $m$ and $n$ are relatively prime positive integers.  Find $m + n.$


 
 \textbf{Answer: }200



\textbf{Solution 1}

The distance between the $x$, $y$, and $z$ coordinates must be even so that the midpoint can have integer coordinates. Therefore, 

 
                \begin{itemize}
                    \item For $x$, we have the possibilities $(0,0)$, $(1,1)$, $(2,2)$, $(0,2)$, and $(2,0)$, $5$ possibilities.\item For $y$, we have the possibilities $(0,0)$, $(1,1)$, $(2,2)$, $(3,3)$, $(0,2)$, $(2,0)$, $(1,3)$, and $(3,1)$, $8$ possibilities.\item For $z$, we have the possibilities $(0,0)$, $(1,1)$, $(2,2)$, $(3,3)$, $(4,4)$, $(0,2)$, $(0,4)$, $(2,0)$, $(4,0)$, $(2,4)$, $(4,2)$, $(1,3)$, and $(3,1)$, $13$ possibilities.
                \end{itemize}


                However, we have $3\cdot 4\cdot 5 = 60$ cases where we have simply taken the same point twice, so we subtract those. Therefore, our answer is $\frac {5\cdot 8\cdot 13 - 60}{60\cdot 59} = \frac {23}{177}\Longrightarrow m+n = \boxed{200}$.

 





\textbf{Solution 2}

There are $(2 + 1)(3 + 1)(4 + 1) = 60$ points in total. We group the points by parity of each individual coordinate -- that is, if $x$ is even or odd, $y$ is even or odd, and $z$ is even or odd. Note that to have something that works, the two points must have this same type of classification (otherwise, if one doesn't match, the resulting sum for the coordinates will be odd at that particular spot).

 There are $12$ EEEs (the first position denotes the parity of $x,$ the second $y,$ and the third $z.$), $8$ EEOs, $12$ EOEs, $6$ OEEs, $8$ EOOs, $4$ OEOs, $6$ OOEs, and $4$ OOOs. Doing a sanity check, $12 + 8 + 12 + 6 + 8 + 4 + 6 + 4 = 60,$ which is the total number of points.

 Now, we can see that there are $12 \cdot 11$ ways to choose two EEEs (respective to order), $8 \cdot 7$ ways to choose two EEOs, and so on. Therefore, we get 
 \[12\cdot11 + 8\cdot7 + 12\cdot11 + 6\cdot5 + 8\cdot7 + 4\cdot3 + 6\cdot5 + 4\cdot3 = 460\] ways to choose two points where order matters. There are $60 \cdot 59$ total ways to do this, so we get a final answer of 
 \[\dfrac{460}{60 \cdot 59} = \dfrac{23}{3 \cdot 59} = \dfrac{23}{177},\] for our answer of $23 + 177 = \boxed{200}.$

 Solution by Ilikeapos

 





\textbf{Solution 3}

Similarly to Solution 2, we note that there are $60$ points and that the parities of the two points' coordinates must be the same in order for the midpoint to be in $S$.

 Ignore the distinct points condition. The probability that the midpoint is in $S$ is then
 \[\left(\left(\frac 23\right)^2+\left(\frac 13\right)^2\right)\left(\left(\frac 24\right)^2+\left(\frac 24\right)^2\right)\left(\left(\frac 35\right)^2+\left(\frac 25\right)^2\right)=\frac{13}{90}.\]

 Note that $\frac{13}{90}=\frac{520}{3600}$. Since there are $3600$ total ways to choose $2$ points from $S$, there must be $520$ pairs of points that have their midpoint in $S$. Of these pairs, $60$ of them contain identical points (not distinct).

 Subtracting these cases, our answer is $\frac{520-60}{3600-60}=\frac{23}{177}\implies\boxed{200}$.

 





\textbf{Solution 4}

There are $(2 + 1)(3 + 1)(4 + 1) = 60$ points in total. Note that in order for the midpoint of the line segment to be a lattice point, the lengths on the x, y, and z axis must be even numbers. We will define all segments by denoting the amount they extend in each dimension: $(x, y, z)$. For example, the longest diagonal possible will be $(2,3,4)$, the space diagonal of the box. Thus, any line segment must have dimensions that are even. For $x$ the segment may have a value of $0$ for $x$, (in which case the segment would be two dimensional) or a value of $2$. The same applies for $y$, because although it is three units long the longest even integer is two. For $z$ the value may be $0$, $2$, or $4$. Notice that if a value is zero, then the segment will pertain to only two dimensions. If two values are zero then the line segment becomes one dimensional. 

 Then the total number of possibilities will be $2 \cdot 2 \cdot 3$.

 Listing them out appears as follows:

 $2,2,4$

 $2,2,2$

 $2,2,0$

 $2,0,4$

 $2,0,2$

 $2,0,0$

 $0,2,4$

 $0,2,2$

 $0,2,0$

 $0,0,4$

 $0,0,2$

 $0,0,0$ * this value is a single point

 
Now, picture every line segment to be the space diagonal of a box. Allow this box to define the space the segment occupies. The question now transforms into "how many ways can we arrange this smaller box in the two by three by four?".

 Notice that the amount an edge can shift inside the larger box is the length of an edge of the larger box (2, 3, or 4) minus the edge of the smaller box (also known as the edge),  plus one. For example, (0, 2, 2) would be $3 \cdot 2 \cdot 3$. Repeat this process.

 
$2,2,4$ 2

 $2,2,2$ 6

 $2,2,0$ 10

 $2,0,4$ 4 

 $2,0,2$ 12

 $2,0,0$ 20

 $0,2,4$ 6

 $0,2,2$ 18

 $0,2,0$ 30

 $0,0,4$ 12

 $0,0,2$ 36

 $0,0,0$ 60 * this won't be included, but notice that sixty the number of lattice points

 
Finally, we remember that there are four distinct space diagonals in a box, so we should multiply every value by four, right? Unfortunately we forgot to consider that some values have only one or two dimensions. They should be multiplied by one or two, respectively. This is because segments with two dimensions are the diagonals of a rectangle and thus have two orientations. Then any value on our list without any zeroes will be multiplied by four, and any value on our list with only one zero will be multiplied by two, and finally any value on our list with two zeroes will be multiplied by one:

 $2,2,4$ 2 8

 $2,2,2$ 6 24

 $2,2,0$ 10 20

 $2,0,4$ 4 8

 $2,0,2$ 12 24

 $2,0,0$ 20 20

 $0,2,4$ 6 12

 $0,2,2$ 18 36

 $0,2,0$ 30 30

 $0,0,4$ 12 12

 $0,0,2$ 36 36

 $0,0,0$ 60 * it's nice to point out that this value will be multiplied by zero

 
add every value on the rightmost side of each term and we will receive $230$. Multiply by two because each segment can be flipped, to receive $460$. There are $60 \cdot 59$ ways to choose two distinct points, so we get 
 \[\dfrac{460}{60 \cdot 59} = \dfrac{23}{3 \cdot 59} = \dfrac{23}{177},\] for our answer of $23 + 177 = \boxed{200}$.

 
Solution by jackshi2006

 




	\item (\textit{AIME I 2001}) In a rectangular array of points, with 5 rows and $N$ columns, the points are numbered consecutively from left to right beginning with the top row.  Thus the top row is numbered 1 through $N,$ the second row is numbered $N + 1$ through $2N,$ and so forth.  Five points, $P_1, P_2, P_3, P_4,$ and $P_5,$ are selected so that each $P_i$ is in row $i.$  Let $x_i$ be the number associated with $P_i.$  Now renumber the array consecutively from top to bottom, beginning with the first column.  Let $y_i$ be the number associated with $P_i$ after the renumbering.  It is found that $x_1 = y_2,$ $x_2 = y_1,$ $x_3 = y_4,$ $x_4 = y_5,$ and $x_5 = y_3.$  Find the smallest possible value of $N.$


 
 \textbf{Answer: }149



\textbf{Solution}

Let each point $P_i$ be in column $c_i$. The numberings for $P_i$ can now be defined as follows.
 \begin{align*}x_i &= (i - 1)N + c_i\\ y_i &= (c_i - 1)5 + i \end{align*}


 We can now convert the five given equalities.
 \begin{align}x_1&=y_2 & \Longrightarrow & & c_1 &= 5 c_2-3\\ x_2&=y_1 & \Longrightarrow & & N+c_2 &= 5 c_1-4\\ x_3&=y_4 & \Longrightarrow & & 2 N+c_3 &= 5 c_4-1\\ x_4&=y_5 & \Longrightarrow & & 3 N+c_4 &= 5 c_5\\ x_5&=y_3 & \Longrightarrow & & 4 N+c_5 &= 5 c_3-2 \end{align}Equations $(1)$ and $(2)$ combine to form
 \[N = 24c_2 - 19\]Similarly equations $(3)$, $(4)$, and $(5)$ combine to form
 \[117N +51 = 124c_3\]Take this equation modulo 31
 \[24N+20\equiv 0 \pmod{31}\]And substitute for N
 \[24 \cdot 24 c_2 - 24 \cdot 19 +20\equiv 0 \pmod{31}\]
 \[18 c_2 \equiv 2 \pmod{31}\]

 Thus the smallest $c_2$ might be is $7$ and by substitution $N = 24 \cdot 7 - 19 = 149$

 \begin{verbatim}
Note: In this problem there are two "groups" of equations that are linked together, that is $(1) + (2)$ and $(3) + (4) + (5)$. For your second congruence used to narrow down $N$, you can use any one of $(3), (4), (5)$. Moreover, you can directly solve for N. The other two equations and the alternate continuation are given below:
\end{verbatim}$73N=124c_4-35$ [combined equation]

 $73N\equiv -35\pmod{124}$ [written mod 124]

 $11N\equiv -4\pmod{31}$ [split the equation mod 4 and mod 31, $N\equiv5\pmod{24}$ supersedes this equation mod 4]

 $2N\equiv 19\pmod{31}$ [multiply both sides by 3, simplify]

 $N\equiv 25\pmod{31}$ [multiply both sides by 16, simplify]

 Then use the \href{https://artofproblemsolving.com/wiki/index.php?title=Chinese\_Remainder\_Theorem}{Chinese Remainder Theorem} with $N\equiv5\pmod{24}$ to solve for N. 

 
$89N=124c_5-74$

 $N\equiv 50\pmod{124}$Continue as above.

 
Check to make sure the values of each $c$ fit the conditions of the problem.The column values can also easily be found by substitution
 \begin{align*}c_1&=32\\ c_2&=7\\ c_3&=141\\ c_4&=88\\ c_5&=107 \end{align*}
As these are all positive and less than $N$, $\boxed{149}$ is the solution.

 \textbf{Sidenote} If we express all the $c_i$ in terms of $N$, we have
 \[24c_1=5N+23\]
 \[24c_2=N+19\]
 \[124c_3=117N+51\]
 \[124c_4=73N+35\]
 \[124c_5=89N+7\]

 It turns out that there exists such an array satisfying the problem conditions if and only if
 \[N\equiv 149 \pmod{744}\]

 


 In addition, the first two equation can be written $n = 5mod24$, and chasing variables in the last three equation gives us $89n + 7 = 124e$. With these two equations you may skip a lot of rewriting and testing. $\boxed{149}$ still appears as our answer.

 
-jackshi2006

 




	\item (\textit{AIME I 2001}) A sphere is inscribed in the tetrahedron whose vertices are $A = (6,0,0), B = (0,4,0), C = (0,0,2),$ and $D = (0,0,0).$  The radius of the sphere is $m/n,$ where $m$ and $n$ are relatively prime positive integers.  Find $m + n.$


 
 \textbf{Answer: }005



\textbf{Solution 1}


                    \begin{figure}[H]
                        \centering
                        \includegraphics[width=0.5\textwidth]{Solution/AIME_I_2001_12_Solution_Figure_1.png}
                    \end{figure}
                    

 The center $I$ of the insphere must be located at $(r,r,r)$ where $r$ is the sphere's radius.$I$ must also be a distance $r$ from the plane $ABC$

 The signed distance between a plane and a point $I$ can be calculated as $\frac{(I-G) \cdot P}{|P|}$, where G is any point on the plane, and P is a vector perpendicular to ABC.

 A vector $P$ perpendicular to plane $ABC$ can be found as $V=(A-C)\times(B-C)=\langle 8, 12, 24 \rangle$

 Thus $\frac{(I-C) \cdot P}{|P|}=-r$ where the negative comes from the fact that we want $I$ to be in the opposite direction of $P$

 
 \begin{align*}\frac{(I-C) \cdot P}{|P|}&=-r\\ \frac{(\langle r, r, r \rangle-\langle 0, 0, 2 \rangle) \cdot P}{|P|}&=-r\\ \frac{\langle r, r, r-2 \rangle \cdot \langle 8, 12, 24 \rangle}{\langle 8, 12, 24 \rangle}&=-r\\ \frac{44r -48}{28}&=-r\\ 44r-48&=-28r\\ 72r&=48\\ r&=\frac{2}{3} \end{align*}


 
Finally $2+3=\boxed{005}$

 





\textbf{Solution 2}

Notice that we can split the tetrahedron into $4$ smaller tetrahedrons such that the height of each tetrahedron is $r$ and the base of each tetrahedron is one of the faces of the original tetrahedron. This is because the bases of the tetrahedrons are tangent to the sphere, so the line from the center to the foot of the perpendicular to the bases hits the tangency points. Letting volume be $V$ and surface area be $F$, using the volume formula for each pyramid(base times height divided by 3) we have $\dfrac{rF}{3}=V$. The surface area of the pyramid is $\dfrac{6\cdot{4}+6\cdot{2}+4\cdot{2}}{2}+[ABC]=22+[ABC]$. We know triangle ABC's side lengths, $\sqrt{2^{2}+4^{2}}, \sqrt{2^{2}+6^{2}},$ and $\sqrt{4^{2}+6^{2}}$, so using the expanded form of heron's formula, 
 \begin{align*}[ABC]&=\sqrt{\dfrac{2(a^{2}b^{2}+b^{2}c^{2}+a^{2}c^{2})-a^{4}-b^{4}-c^{4}}{16}}\\ &=\sqrt{2(5\cdot{13}+10\cdot{5}+13\cdot{10})-5^{2}-10^{2}-13^{2}}\\ &=\sqrt{196}\\ &=14\end{align*}
Therefore, the surface area is $14+22=36$, and the volume is $\dfrac{[BCD]\cdot{6}}{3}=\dfrac{4\cdot{2}\cdot{6}}{3\cdot{2}}=8$, and using the formula above that $\dfrac{rF}{3}=V$, we have $12r=8$ and thus $r=\dfrac{2}{3}$, so the desired answer is $2+3=\boxed{005}$.

 (Solution by Shaddoll)

 





\textbf{Solution 3}

The intercept form equation of the plane $ABC$ is $\frac{x}{6}+\dfrac{y}{4}+\dfrac{z}{2}=1.$ Its normal form is $\dfrac{2}{7}x+\dfrac{3}{7}y+\dfrac{6}{7}z-\dfrac{12}{7}=0$ (square sum of the coefficients equals 1). The distance from $(r,r,r)$ to the plane is $\left |\dfrac{2}{7}r+\dfrac{3}{7}r+\dfrac{6}{7}r-\dfrac{12}{7}\right |$. Since $(r,r,r)$ and $(0,0,0)$ are on the same side of plane, the value in the absolute value sign is negative (same as the one by plugging in $(0,0,0)$). Therefore we have $-\left (\dfrac{2}{7}r+\dfrac{3}{7}r+\dfrac{6}{7}r-\dfrac{12}{7}\right )=r.$ So $r=\dfrac{2}{3},$ which solves the problem.

 Additionally, if $(r,r,r)$ is on the other side of $ABC$, we have $\left (\dfrac{2}{7}r+\dfrac{3}{7}r+\dfrac{6}{7}r-\dfrac{12}{7}\right )=r$, which yields $r=\dfrac{12}{5},$ corresponding an "ex-sphere" that is tangent to face $ABC$ as well as the extensions of the other 3 faces.

 -JZ

 





\textbf{Solution 4}

First let us find the equation of the plane passing through $(6,0,0), (0,0,2), (0,4,0)$.  The "point-slope form" is $A(6-x1)+B(0-y1)+C(0-z1)=0.$ Plugging in $(0,0,2)$ gives $A(6)+B(0)+C(-2)=0.$  Plugging in $(0,4,0)$ gives $A(6)+B(-4)+C(0)=0.$  We can then use Cramer's rule/cross multiplication to get $A/(0-8)=-B/(0+12)=C/(-24)=k.$  Solve for A, B, C to get $2k, 3k, 6k$ respectively.  We can then get $2k(x-x1)+3k(y-y1)+6k(z-z1)=0.$  Cancel out k on both sides. Next, let us substitute $(0,0,2)$.  We can then get $2x+3y+6z=12$as the equation of the plane.  We can divide the equation by its magnitude to get the normal form of the plane.  We get $2x/7+2y/7+6z/7=12/7$ to be the normal form.  Note that the point is going to be at $(r,r,r).$  We find the distance from $(r,r,r)$ to the plane as $2/7r+3/7r+6/7r-12/7/(\sqrt{(4/49+9/49+36/49)})$, which is $+/-(11r/7-12/7)$.  We take the negative value of this because if we plug in $(0,0,0)$ to the equation of the plane we get a negative value.  We equate that value to r and we get the equation $-(11r/7-12/7)=r$ to solve $r={2/3}$, so the answer is $\boxed{005}$.

 





\textbf{Solution 5}

Clearly, if the radius of the sphere is $r$, the center of the sphere lies on $(r, r, r)$.

 We find the equation of plane $ABC$ to be $\frac16 x+\frac14 y+\frac12 z=1$. From the definition of the insphere, it must be true that the distance from the center of the sphere to plane $ABC$ is equal to the length of the radius of the sphere. By point-to-plane, we have 
 \[r=\frac{|\frac16 r+\frac14 r+\frac12 r-1|}{\sqrt{\left(\frac16\right)^2+\left(\frac14\right)^2+\left(\frac12\right)^2}} \implies r=\frac23,\] so the answer is $\boxed{005}$.

 -pqr.

 





\textbf{Solution 6(Formula Bash)}

The radius of the insphere in a tetrahedron can be calculated using the formula $r = \frac{3V}{A}$, where $r$ is the radius, $V$ is the volume of the tetrahedron, and $A$ is the total surface area of the tetrahedron.

 We calculate the volume of the tetrahedron as:
 \[V = \frac{2 \times 4 \times 6}{6} = 8\]

 Next, we find the total surface area as the sum of the areas of the four triangular faces:
 \[\text{Surface Area} = [ABD] + [BDC] + [ACD] + [ABC] = 12 + 4 + 6 + [ABC]\]

 The side lengths of triangle $ABC$ are $\sqrt{52}$, $\sqrt{40}$, and $\sqrt{20}$. Constructing an altitude from $A$ and using the system of equations $x^2 + h^2 = 40$ and $(2\sqrt{5} - x)^2 + h^2 = 20$, we solve for $h$ and get:
 \[h = \frac{14\sqrt{5}}{5}\]

 Thus, the area of triangle $ABC$ is:
 \[\text{Area of } [ABC] = 14\]

 Now, the total surface area is:
 \[A = 12 + 4 + 6 + 14 = 36\]Finally, using the formula for the radius, we have:
 \[r = \frac{3 \times 8}{36} = \frac{2}{3}\]$\boxed{005}$

 




	\item (\textit{AIME I 2001}) In a certain circle, the chord of a $d$-degree arc is 22 centimeters long, and the chord of a $2d$-degree arc is 20 centimeters longer than the chord of a $3d$-degree arc, where $d < 120.$  The length of the chord of a $3d$-degree arc is $- m + \sqrt {n}$ centimeters, where $m$ and $n$ are positive integers.  Find $m + n.$


 
 \textbf{Answer: }174



\textbf{Solution 1}


                    \begin{figure}[H]
                        \centering
                        \includegraphics[width=0.5\textwidth]{Solution/AIME_I_2001_13_Solution_Figure_1.png}
                    \end{figure}
                    Note that a cyclic quadrilateral in the form of an isosceles trapezoid can be formed from three chords of three $d$-degree arcs and one chord of one $3d$-degree arc. The diagonals of this trapezoid turn out to be two chords of two $2d$-degree arcs. Let $AB$, $AC$, and $BD$ be the chords of the $d$-degree arcs, and let $CD$ be the chord of the $3d$-degree arc. Also let $x$ be equal to the chord length of the $3d$-degree arc. Hence, the length of the chords, $AD$ and $BC$, of the $2d$-degree arcs can be represented as $x + 20$, as given in the problem. 

 Using \href{https://artofproblemsolving.com/wiki/index.php?title=Ptolemy\%27s\_theorem}{Ptolemy's theorem},

 
 \[AB \cdot CD + AC \cdot BD = AD \cdot BC\]
 \[\iff 22x + 22 \cdot 22  = (x + 20)^2\]
 \[\iff 22x + 484 = x^2 + 40x + 400\]
 \[\iff x^2 + 18x - 84 = 0.\]

 We can then apply the quadratic formula to find the positive root of this equation (since polygons obviously cannot have sides of negative length):
 \[x = \frac{-18 + \sqrt{18^2 + 4 \cdot 84}}{2} = \frac{-18 + \sqrt{660}}{2}.\]

 This simplifies to $x = \frac{-18 + 2\sqrt{165}}{2} = -9 + \sqrt{165}$. Thus the answer is $9 + 165 = \boxed{174}$.

 





\textbf{Solution 2}

Let $z=\frac{d}{2},$ and $R$ be the circumradius. From the given information, 
 \[2R\sin z=22, \quad\text{and}\] 
 \[2R(\sin 2z-\sin 3z)=20.\] Dividing the latter equation by the former gives 
 \[\frac{2\sin z\cos z-\left(3\cos^2z\sin z-\sin^3 z\right)}{\sin z}=2\cos z-\left(3\cos^2z-\sin^2z\right)=1+2\cos z-4\cos^2z=\frac{10}{11}\] 
 \[\iff 4\cos^2z-2\cos z-\frac{1}{11}=0. \qquad (*)\] We want to find 
 \[\frac{22\sin (3z)}{\sin z}=22(3-4\sin^2z)=22(4\cos^2z-1).\] From $(*),$ this is equivalent to $44\cos z-20$. Using the quadratic formula, we deduce that this expression equals $-9+\sqrt{165}$, so our answer is $\boxed{174}$.

 





\textbf{Solution 3}

Let $z=\frac{d}{2}$, $R$ be the circumradius, and $a$ be the length of a $3d$-degree chord. Using the extended sine law, we obtain:
 \[22=2R\sin(z),\]
 \[20+a=2R\sin(2z), \quad\text{and}\]
 \[a=2R\sin(3z).\]Dividing the second equation by the first, and using the double angle formula, we obtain $\cos(z)=\frac{20+a}{44}$.Now, using the triple angle formula, we can rewrite the third equation as follows: 
 \[a=2R \sin(3z)=\frac{22}{\sin(z)}\left(3\sin(z)-4\sin^3(z)\right) = 22\left(3-4\sin^2(z)\right) = 22\left(4\cos^2(z)-1\right) = \frac{(20+a)^2}{22}-22,\]and solving this quadratic equation gives the answer as $\boxed{174}$.

 




	\item (\textit{AIME I 2001}) A mail carrier delivers mail to the nineteen houses on the east side of Elm Street.  The carrier notices that no two adjacent houses ever get mail on the same day, but that there are never more than two houses in a row that get no mail on the same day.  How many different patterns of mail delivery are possible?


 
 \textbf{Answer: }351



\textbf{Solution 1}

Let $0$ represent a house that does not receive mail and $1$ represent a house that does receive mail. This problem is now asking for the number of $19$-digit strings of $0$'s and $1$'s such that there are no two consecutive $1$'s and no three consecutive $0$'s. 

 The last two digits of any $n$-digit string can't be $11$, so the only possibilities are $00$, $01$, and $10$. 

 Let $a_n$ be the number of $n$-digit strings ending in $00$, $b_n$ be the number of $n$-digit strings ending in $01$, and $c_n$ be the number of $n$-digit strings ending in $10$. 

 If an $n$-digit string ends in $00$, then the previous digit must be a $1$, and the last two digits of the $n-1$ digits substring will be $10$. So 
 \[a_{n} = c_{n-1}.\]

 If an $n$-digit string ends in $01$, then the previous digit can be either a $0$ or a $1$, and the last two digits of the $n-1$ digits substring can be either $00$ or $10$. So 
 \[b_{n} = a_{n-1} + c_{n-1}.\]

 If an $n$-digit string ends in $10$, then the previous digit must be a $0$, and the last two digits of the $n-1$ digits substring will be $01$. So 
 \[c_{n} = b_{n-1}.\]

 Clearly, $a_2=b_2=c_2=1$. Using the \href{https://artofproblemsolving.com/wiki/index.php?title=Recursion}{recursive} equations and initial values: 
 \[\begin{array}{|c|c|c|c|c|c|c|c|c|c|c|c|c|c|c|c|c|c|c|} \multicolumn{19}{c}{}\\\hline n&2&3&4&5&6&7&8&9&10&11&12&13&14&15&16&17&18&19\\\hline a_n&1&1&1&2&2&3&4&5&7&9&12&16&21&28&37&49&65&86\\\hline b_n&1&2&2&3&4&5&7&9&12&16&21&28&37&49&65&86&114&151\\\hline c_n&1&1&2&2&3&4&5&7&9&12&16&21&28&37&49&65&86&114\\\hline \end{array}\]

 As a result $a_{19}+b_{19}+c_{19}=\boxed{351}$.

 





\textbf{Solution 2 ( Less recursion than solution 1)}

Let $M_n$ represent the number of mail delivery patterns that end with the last house receiving mail. This is $b_n$ in Solution 1. Similarly define $A_n$ to be the number of mail delivery patterns that end with last house not receiving mail. This is just $a_n$ and $c_n$ in solution 1. Let $T_n$ be the total number of mail delivery patterns.

 Here are the possible ending cases: the string ends in $1, 10,$ or $100$. The first case is just $M_n$. The second case is $M_{n-1}$. The third case is $M_{n-2}$. So we have $T_n = M_n + M_{n-1} + M_{n-2}$. Since we want $T_{19}$, it is just $M_{18} + M_{17} + M_{16}$.Now using the same logic as above we can find $M_n = M_{n-2} + M_{n-3}$ ( the cases are 01 and 001). We can refer back to solution 1's table and only keep track of $b_n$, ignoring both $a_n$ and $c_n$. 

 - MathLegend27

 


 





\textbf{Solution 3 (Tedious Casework)}

We split the problem into cases using the number of houses that get mail. Let "|" represent a house that gets mail, and "o" represent a house that doesn't. With a fixed number of |, an o can be inserted between 2 |'s or on the very left or right. There cannot be more than one o that is free to arrange to be placed between two |'s because no three o's can be adjacent, but there can be a maximum of two o's placed on the very left or right. Note that according to the \href{https://artofproblemsolving.com/wiki/index.php?title=Pigeonhole\_Principle}{Pigeonhole Principle}, no more than 10 houses can get mail on the same day. 

 Case 1: 10 houses get mail. No 2 adjacent houses can get mail on the same day, so there must be an o between every two |. $10-1=9$ o's are fixed so we count the number of ways to insert $19 - 10 - 9 = 0$ o's to $10+1 = 11$ spots, or $\binom{11}{0} = 1$.

 Case 2: 9 houses get mail. In this case, $9-1 = 8$ o's are fixed so we count the number of ways to insert $19 - 9 - 8 = 2$ o's to $9+1=10$ spots. However, there is also the case where two o's are both on the very left / right. When both o's that are free to arrange are put on a side, there are $10-1=9$ spots left to insert $2-2=0$ o's. Hence the total number of ways in this case is $\binom{10}{2} + 2\binom{9}{0} = 47$.

 Case 3: 8 houses get mail. In this case, $8-1=7$ o's are fixed so we count the number of ways to insert $19-8-7=4$ o's to $8+1=9$ spots. When two o's are put to the very left / right, there are $9-1=8$ spots left to insert $4-2=2$ o's. We also need to take care of the case where two o's are on the very left and two o's are on the very right: we have $9-1-1=7$ spots to insert $4-2-2=0$ o's. Hence the total number of ways in this case is $\binom{9}{4} + 2\binom{8}{2} + \binom{7}{0} = 183$.

 Case 4: 7 houses get mail. In this case, $7-1=6$ o's are fixed so we count the number of ways to insert $19-7-6=6$ o's to $7+1=8$ spots. When two o's are put to the very left / right, there are $8-1=7$ spots left to insert $6-2=4$ o's. When two o's are on the very left and two o's are on the very right, we have $8-1-1=6$ spots to insert $6-2-2=2$ o's. Hence the total number of ways in this case is $\binom{8}{6} + 2\binom{7}{4} + \binom{6}{2} = 113$.

 Case 5: 6 houses get mail. We have to be careful in this case: $6-1=5$ o's are fixed so we are inserting $19-6-5=8$ o's to $6+1=7$ spots, which means that at least 1 of the 2 sides must have two o's. When 1 of the 2 sides have two o's, there are $7-1=6$ spots to insert $8-2=6$ o's. When both sides have two o's, there are $7-1-1=5$ spots to insert $8-2-2=4$ o's. Hence the total number of ways in this case is $2\binom{6}{6} + \binom{5}{4} = 7$. 

 When less than 6 houses get(s) mail, it's again not possible since at least three o's must be together (again, according to the Pigeonhole Principle). Therefore, the desired answer is $1+47+183+113+7=\boxed{351}$.

 - MP8148

 





\textbf{Solution 4 (Clean Recursion without Bash)}

There doesn't seem to be anything especially noticeable about the number nineteen in this problem, meaning that we can replace the number nineteen with any number without a big effect on the logic that we use to solve the problem. This pits the problem as a likely candidate for recursion.

 At first, it's not immediately clear how to relate the state of $n$ houses in general to that of $n - 1, n - 2,$ or $n - 3.$ We thus break it up into cases, based on whether the first house gets mail or not.

 Let $p_n$ be the number of ways to distribute the mail to $n$ houses. Assume that the first house gets mail. Therefore, since no two adjacent houses get mail on the same day, the second house must not get mail. Starting from the third house, however, things start to look messy, and it looks like we have to break our recurrence down into even smaller cases, which is something that we don't like -- we want to keep our relations as simple as possible. Therefore, seeing that we can't work forwards anymore, we try to work backwards.

 Once the mail carrier delivers the mail to the first and (lack of mail) to the second houses, have him deliver mail to the remaining $n - 3$ houses at the end of the row, skipping the third house. There are $p_{n - 3}$ ways to do this. Now, we see that the availability of mail at the third house is fixed -- if the fourth house doesn't receive mail, the third one must, and if the fourth house receives mail, the third one can't. Therefore, there are simply $p_{n-3}$ ways to deliver the mail if the first house gets mail.

 If the first house doesn't get mail, then we use the same logic -- have the mail carrier skip the second house and deliver the remaining mail to the $n - 2$ houses in $p_{n-2}$ ways. Then, the availability of mail for the second house is fixed, so there are $p_{n - 2}$ ways to deliver the mail in this case.

 We thus have established a recurrence relation -- since the first house either gets mail or it doesn't, and cannot achieve both at the same time, we are confident about the validity of our relation: 
 \[p_n = p_{n-2} + p_{n-3}.\]Now, we simply calculate $p_1, p_2,$ and $p_3.$ Then, it's off to the races for computation!

 $p_1 = 2,$ because the first house can either gets mail or it doesn't -- there are no restrictions.

 $p_2 = 3,$ because all of the possible deliveries are valid (of which there are $2 \cdot 2 = 4$) except the one where both houses receive mail.

 $p_3 = 4,$ as there are $4$ possible ways (here, M represents that that house gets mail and N represents no mail): MNM, MNN, NNM, NMN.

 Using our recurrence relation, we eventually get that $p_{19} = \boxed{351},$ and we're done. 

 Solution by Ilikeapos

 





\textbf{Solution 5}

We use similar wording as in solution 1. In this problem, we divide into 3 cases:

 Case 1: The first house gets mail. In other words, the sequence starts with a $1.$

 We first introduce two variables. Let $x$ be the number of 01's, and let $y$ be the number of 001's in the sequence. For each case, we will divide further into three sub-cases, based on the pattern of mail delivery for the last three houses.

 Sub-case 1: The last house that gets mail is the 17th house. Thus, we have the equation $2x+3y=16,$ which has solutions $(2,4),(5,2),$ and $(8,0).$ There are $\binom{2+4}{2}+\binom{5+2}{2}+\binom{8+0}{8}=37$ total sequences.

 Sub-case 2: The last house that gets mail is the 18th house. We have $2x+3y=17,$ and finding all solutions yields 49 total sequences.

 Sub-case 3: The last house that gets mail is the 19th house. We have $2x+3y=18,$ which gives us 65 total sequences.

 There are 151 sequences for this case.

 Case 2: The first house that gets mail is the 2nd house. The sequence begins with $01.$

 Our sub-cases are still the same. However, our equations become $2x+3y=15,16,17.$ Computing yields $28+37+49=114$ sequences.

 Case 3: The first house that gets mail is the 3rd house. The sequence starts with $001.$

 We have the equations $2x+3y=14,15,16$ so we get $21+28+37=86$ sequences.

 Totaling up gives us $151+114+86=351$ different sequences.

 -math129

 





\textbf{Solution 6 (Different recursion with two rules)}

Let $w_n$ be the number of possible ways if the last house has mail, and $b_n$ be the number of possible ways if the last house does not have mail.

 If the last house has mail, then, the next house can't have mail, meaning that $b_n = w_{n - 1}$. 

 If the last house doesn't have mail, then the next house can either have mail or not have mail. If the next house has mail, then we simply count the number of ways that the row ends in a house with mail, so that means so far, our recursive rule is $w_n = b_{n - 1} + \text{something}$. If the next house does not have mail, then the next house after that must have mail, meaning that $w_n = b_{n - 1} + b_{n - 2}$. 

 Recursing all the way up to $b_{19}$ and $w_{19}$, we get $100 + 251 = \boxed{351}$

 





\textbf{Solution 7 (Simple Recursion)}

Let $a_n$ be the number of ways if the first house has mail, and let $b_n$ be the number of ways if the first house does not get mail.

 $a_n=a_{n-2}+a_{n-3}$ because if the first house gets mail, the next house that gets mail must either be the third or fourth house.

 $b_n=a_{n-1}+a_{n-2}$ because if the first house does not get mail, the next house that gets mail must either be the second or third house.

 Note that we only need list out values of $a_n$ as $b$ depends on $a$. $a_1=1, a_2=1, a_3=2, a_4=2, \ldots$.

 $a_{19}+b_{19}=a_{17}+a_{16}+a_{18}+a_{17}=a_{16}+2\cdot a_{17}+a_{18}=65+2\cdot 86+114=\boxed{351}$.

 





\textbf{Solution 8 (- Correspondence & Constructive Counting)}

Like the previous solutions we use $1$s and $0$s to indicate whether a house get mails or not. Call a sequence of $1$s and $0$s valid if it has no two consecutive $1$s or three consecutive $0$s.

 We claim there is a $1$-to-$1$ correspondence between (A) a valid sequence of length $19$; and (B) a valid sequence of length $25$ that starts and ends with $1$.

 For each valid sequence in (B), we can remove the first three numbers as well as the last three numbers to get a valid sequence in (A). 

 On the other hand, for each valid sequence in (A), we can append three numbers to each end by the following rules to get a valid sequence in (B):

 1) If the sequence starts/ends with $1$, we append $100$/$001$ to the start/end;

 2) If the sequence starts/ends with $0$, we append $101$ to the start/end.

 Now, since all valid sequences in (B) start with either $1010$ or $1001$, and end with either $0101$ or $1001$, it's easy to see this is a $1$-$1$ correspondence.

 To compute the number of valid sequences in (B), we start with $n$ of $1$s with $n-1$ of $0$s in between, then distribute the rest $25-(2n-1)=26-2n$ of $0$s to the $n-1$ slots. The total number of valid sequences is then

 
 \[\sum_n \binom{n-1}{26-2n} = \binom{12}{0} + \binom{11}{2} + \binom{10}{4} + \binom{9}{6} + \binom{8}{8} = 1 + 55 + 210 + 84 + 1 = 351.\]

 





\textbf{Solution 9}

When we see this problem, we immediately think of recursion. Let $a_k$ be the point such that the rightmost house of k houses receives mail. Then $a_1=1$, $a_2=1$, $a_3=2$. Now, we find the recurrence relation 
 \[a_n=a_{n-3}+a_{n-2}\] because  we can either have 1 space in between the mailed rightmost house of $a_n$, or two spaces, as in our previous a’s we already have the a\_n mailed rightmost. Solving, we get $a_{17}=86, a_{18}=114, a_{19}=151$. Now, we realize that the number of ways for the 19 houses to get mail is the sum of $a_{17},a_{18},a_{19}$ because we can either have 2 empty no mail houses, 1 house, or the rightmost house is lit. Summing, we get $\boxed{351}$. 

 




	\item (\textit{AIME I 2001}) The numbers $1, 2, 3, 4, 5, 6, 7,$ and $8$ are randomly written on the faces of a regular octahedron so that each face contains a different number.  The probability that no two consecutive numbers, where $8$ and $1$ are considered to be consecutive, are written on faces that share an edge is $m/n,$ where $m$ and $n$ are relatively prime positive integers.  Find $m + n.$


 
 \textbf{Answer: }085



\textbf{Solution 1}

Choose one face of the octahedron randomly and label it with $1$.  There are three faces adjacent to this one, which we will call A-faces.  There are three faces adjacent to two of the A-faces, which we will call B-faces, and one face adjacent to the three B-faces, which we will call the C-face.

 Clearly, the labels for the A-faces must come from the set $\{3,4,5,6,7\}$, since these faces are all adjacent to $1$.  There are thus $5 \cdot 4 \cdot 3 = 60$ ways to assign the labels for the A-faces.  

 The labels for the B-faces and C-face are the two remaining numbers from the above set, plus $2$ and $8$.  The number on the C-face must not be consecutive to any of the numbers on the B-faces.  

 From here it is easiest to brute force the $10$ possibilities for the $4$ numbers on the B and C faces:

 
                \begin{itemize}
                    \item 2348 (2678): 8(2) is the only one not adjacent to any of the others, so it goes on the C-face.  4(6) has only one B-face it can go to, while 2 and 3 (7 and 8) can be assigned randomly to the last two.  2 possibilities here.
                \end{itemize}


                
                \begin{itemize}
                    \item 2358 (2578): 5 cannot go on any of the B-faces, so it must be on the C-face.  3 and 8 (2 and 7) have only one allowable B-face, so just 1 possibility here.
                \end{itemize}


                
                \begin{itemize}
                    \item 2368 (2478): 6(4) cannot go on any of the B-faces, so it must be on the C-face.  3 and 8 (2 and 7) have only one allowable B-face, so 1 possibility here.
                \end{itemize}


                
                \begin{itemize}
                    \item 2458 (2568): All of the numbers have only one B-face they could go to.  2 and 4 (6 and 8) can go on the same, so one must go to the C-face.  Only 2(8) is not consecutive with any of the others, so it goes on the C-face.  1 possibility.
                \end{itemize}


                
                \begin{itemize}
                    \item 2378: None of the numbers can go on the C-face because they will be consecutive with one of the B-face numbers.  So this possibility is impossible.
                \end{itemize}


                
                \begin{itemize}
                    \item 2468: Both 4 and 6 cannot go on any B-face.  They cannot both go on the C-face, so this possibility is impossible.
                \end{itemize}


                There is a total of $10$ possibilities.  There are $3!=6$ permutations (rotations/reflections) of each, so $60$ acceptable ways to fill in the rest of the octahedron given the $1$.  There are $7!=5040$ ways to randomly fill in the rest of the octahedron.  So the probability is $\frac {60}{5040} = \frac {1}{84}$. The answer is $\boxed{085}$.

 





\textbf{Solution 2}

Consider the cube formed from the face centers of the regular octahedron. Color the vertices in a checker board fashion. We seek the number of circuits traversing the cube entirely composed of diagonals. Notice for any vertex, it can be linked to at most one different-colored vertex, i.e. its opposite vertex. Thus, there are only two possible types of circuits ($B$ for black and $W$ for white).

 Type I: $BB-WWWW-BB$. There are $4!$ ways to arrange the black vertices and consequently two of the white vertices and $2!$ ways to arrange the other two white vertices. Since the template has a period of $8$, there are $4!\cdot 2!\cdot 8 = 384$ circuits of type I.

 Type II: $B-WW-BB-WW-B$. There are $4!$ ways to arrange the black vertices and consequently the white vertices. Since the template has a period of $4$, there are $4! \cdot 4 = 96$ circuits of type II.

 Thus, there are $384+96=480$ circuits satisfying the given condition, out of the $8!$ possible circuits. Therefore, the desired probability is $\frac{480}{8!} = \frac{1}{84}$. The answer is $\boxed{085}$.

 





\textbf{Solution 3}

As in the previous solution, consider the cube formed by taking each face of the octahedron as a vertex. Let one fixed vertex be A. Then each configuration (letting each vertex have a number value from 1-8) of A and the three vertices adjacent to A uniquely determine a configuration that satisfies the conditions, i.e. no two vertices have consecutive numbers. Thus, the number of desired configurations is equivalent to the number of ways of choosing the values of A and its three adjacent vertices. 

 The value of A can be chosen in 8 ways, and the 3 vertices adjacent to A can be chosen in $5\cdot4\cdot3=60$ ways, since they aren't adjacent to each other, but they can't, after all, be consecutive values to A. For example, if A=1, then the next vertex can't be 1,2 or 8, so there are 5 choices. However, the next vertex also adjacent to A can be chosen in 4 ways; it can't be equal to 1,2,8, or the value of the previously chosen vertex. With the same reasoning, the last such vertex has 3 possible choices.

 The total number of ways to choose the values of the vertices of the cube independently is 8!, so our probability is thus $\frac{8\cdot60}{8!}=\frac{8\cdot5\cdot4\cdot3}{8!}=\frac{1}{84}$, from which the answer is $\boxed{085}$.

 , the average is one per choice of 3 to border the 1, ex 5,6,7 border 2 solutions, 3,5,7 border 0 solution

 





\textbf{Solution 4}


                    \begin{figure}[H]
                        \centering
                        \includegraphics[width=0.5\textwidth]{Solution/AIME_I_2001_15_Solution_Figure_1.png}
                    \end{figure}
                    

 The probability is equivalent to counting the number of Hamiltonian cycles in this 3D graph over $7!.$ This is because each Hamiltonian cycle corresponds to eight unique ways to label the faces. Label the vertices $AR,BR,CR,DR,AX,BX,CX,DX$ where vertices $ab$ and $cd$ are connected if $a=c$ or $b=d.$

 Case 1: Four of the vertical edges are used. $6\cdot 2=12.$

 Case 2: Two of the vertical edges are used. $4\cdot 3 \cdot 2\cdot 2=48.$

 So, the probability is $\frac{60}{5040}=\frac{1}{84}.$ Therefore, our answer is $\boxed{085}$

 





\textbf{Solution 5}

As with some of the previous solutions, consider the cube formed by connecting the centroids of the faces on the octahedron. We choose a random vertex(hence fixing the diagram), giving us $7!$ ways as our denominator. WLOG, we color this start vertex red, and we color all $3$ vertices adjacent to it blue. We repeat this for the other vertices.

 Note that there is a 1-1 correspondence between the number of valid face numberings with rotational symmetry and the number of ways to move a particle to every vertex of the cube and returning to the start vertex with only diagonal moves. We can move along either short diagonals and long diagonals.

 Next, note that we can only move from the red vertices to the blue vertices and back with long diagonals(since short diagonals keep us on the color we are currently on). Thus, it is easy to check that the only possible sequences of long and short diagonals are cyclic permutations of $SSSLSSSL$ and $SLSLSLSL$.

 Case 1: $SSSLSSSL$.There are $4$ cyclic permutations, and we can traverse the entirety of the red tetrahedron in $4$ steps in $3!$ ways. Then, after the move to the blue tetrahedron from one of the non-starting red vertices, we realize that our first step cannot be to the blue vertex opposite the starting red vertex. Hence, there are $2$ possibilities for our first step. However, once we make our first move, the path is fixed. This leaves us with $4! \cdot 2$ ways.

 Case 2: $SLSLSLSL$There are $2$ cyclic permutations in this case. After we choose one of the $3$ red vertices to go to from the starting vertex and move to the blue tetrahedron, we are once again left with $2$ choices to move to by similar logic to the first case. However, after we move back to the red tetrahedron, our choices are fixed. This leaves us with $2 \cdot 3 \cdot 2$ ways.

 Finally, summing and dividing by $7!$ gives us $\frac{60}{7!} = \frac{1}{84} \rightarrow \boxed{085}$, as desired. - Spacesam

 


 


 


 




\end{enumerate}